% =========================================================================
\section{Conclusion and Research Roadmap}
\label{sec:conclusion}

We presented FL-Iroh, a federated learning system for agricultural and environmental IoT that embeds NAT traversal in the application transport, addresses and admits peers by public key, and coordinates them through a lightweight CoAP control plane with a resource directory. On real hardware across five networks with RFC~4787-classified NAT behavior, including a mobile-carrier CGNAT, the system kept clients and aggregator mutually reachable without static IPs, port forwarding or VPN overlays: long-lived sessions used validated direct paths, the encrypted relay delivered every transfer when UDP was blocked, and transfers resumed after link outages and network switches. A complete federation over the WAN ran on a Raspberry Pi~3B client at about one second of training and 0.4~GB of memory per round, all rounds of a loss/latency/bandwidth sweep completed with unchanged accuracy for edge-sized models, unauthorized peers were rejected before any payload was read, and resource-directory discovery remained in the tens of milliseconds for 5{,}000 endpoints. The evaluation also delimits the system: cold-start connections are often relay-carried during their first seconds, direct paths through mobile CGNAT can be less reliable than the relay, the Python/PyTorch runtime is heavy for the smallest devices, and partial aggregation for GAMs brings no benefit when local history is long. FL-Iroh keeps raw data on the device and encrypts updates in transit; it is not, by itself, a privacy-preserving learning system.

Table~\ref{tab:roadmap} organizes the open problems into a research roadmap, from engineering steps that build directly on the released artifact to open research questions for the community.

\begin{table*}[t]
\centering
\caption{Research roadmap for NAT-transparent federated learning in agricultural and environmental IoT}
\label{tab:roadmap}
\footnotesize
\renewcommand{\arraystretch}{1.2}
\begin{tabular}{p{2.4cm} p{5.6cm} p{7.4cm}}
\toprule
\textbf{Horizon} & \textbf{Direction} & \textbf{Concrete opportunities} \\
\midrule
Short term (engineering) & Measurement breadth & Multi-operator, multi-country connectivity campaigns with RFC~4787 probes; many clients behind one CGNAT; IPv6-only networks; native (non-WSL) aggregators; field validation of endpoint recovery. \\
 & Operational baselines & Comparisons with other overlays and frameworks across sites and on larger fleets, extending the single-network live comparison with Flower over Tailscale; self-hosted relays and their capacity. \\
 & Measured energy & Shunt-monitor or PMIC-based energy per round on solar-powered nodes; energy-aware client selection using the capability tags of the resource directory. \\
\midrule
Medium term (systems) & Lightweight runtimes & Framework-neutral tensor encoding; native Rust/C clients for MCU-class devices; integer and sparse update formats. \\
 & Communication efficiency & Update compression and quantization tuned to relay versus direct paths; path-aware scheduling (e.g., large transfers over the relay on mobile CGNAT). \\
 & Asynchrony and intermittency & Asynchronous and buffered aggregation for multi-day disconnections; server-driven termination and rejoin semantics. \\
\midrule
Long term (research) & Trustworthy FL & Secure aggregation and differential privacy over P2P transports; robust aggregation against poisoned updates; key revocation and rotation for federations. \\
 & Decentralized topologies & Gossip or hierarchical FL over public-key addressed peers without a central aggregator; relay placement as an optimization problem. \\
 & Domain models & Federated time-series models (GAMs, state-space models) with short local histories, where partial aggregation may matter; concept drift in multi-year sensor data. \\
\bottomrule
\end{tabular}
\end{table*}
